% Please do not change %
\documentclass[11pt]{article} % font size
\usepackage[margin=1in]{geometry} % margins
\usepackage{amsmath,amsthm,amssymb} % for the  common "math symbols"
\usepackage{algorithm}
\usepackage{algpseudocode}
\usepackage[shortlabels]{enumitem} % for enumeration
\usepackage{booktabs} % if you need tables
\usepackage{listings}
\usepackage{color}
\usepackage{quiver}
\usepackage{ulem}
\DeclareMathOperator*{\argmax}{argmax}
\DeclareMathOperator*{\argmin}{argmin}

% \theoremstyle{definition}
\newtheorem{theorem}{Theorem}[section]
\newtheorem{corollary}[theorem]{Corollary}
\newtheorem{lemma}[theorem]{Lemma}
\newtheorem{proposition}[theorem]{Proposition}
\newtheorem{conjecture}[theorem]{Conjecture}
\theoremstyle{definition}
\newtheorem{definition}[theorem]{Definition}
\theoremstyle{definition}
\newtheorem{fact}[theorem]{Fact}
\theoremstyle{definition}
\newtheorem{example}[theorem]{Example}

\lstset{
  frame=none,
  xleftmargin=2pt,
  stepnumber=1,
  numbers=left,
  numbersep=5pt,
  numberstyle=\ttfamily\tiny\color[gray]{0.3},
  belowcaptionskip=\bigskipamount,
  captionpos=b,
  escapeinside={*'}{'*},
  language=haskell,
  tabsize=2,
  emphstyle={\bf},
  commentstyle=\it,
  stringstyle=\mdseries\rmfamily,
  showspaces=false,
  keywordstyle=\bfseries\rmfamily,
  columns=flexible,
  basicstyle=\small\sffamily,
  showstringspaces=false,
  morecomment=[l]\%,
}
% Feel free to include additional packages (if needed), but make sure it does not override the margin/font requirements.
%
\setlength{\parindent}{0pt} % no indent for new paragraphs.
\setlength{\parskip}{5pt} % distance between paragraphs, which will be used when you have a blank line in your LaTeX code.


%%%%%%%%%%%%%%%%%%%%%%%%%%%%%%%%%%%%%%%%%%%%%%%%%%%%%%%%%%%%%%%%%%%%%%%%%%%%%%
\title{Categories, Proofs, and Processes: \\Assignment 4}
%
\author{Wira Azmoon Ahmad}
\date{19 Feb 2024}
%
\begin{document}
%
\maketitle

%You may use the usual \LaTeX\ environments to structure your answers:
%\begin{align}
%        \label{basegnn}
%        \mathbf{h}_u^{(t+1)} = \sigma \bigg( \mathbf{W}_\text{self}^{(t)} \mathbf{h}_u^{(t)}  + \mathbf{W}_\text{neigh}^{(t)} \sum_{v \in N(u)} \mathbf{h}_v^{(t)} + \mathbf{b}^{(t)} \bigg): 
%\end{align}


%\begin{table}[h]
%\caption{A simple table.}
%\centering
%\begin{tabular}[t]{lrrrr}
%\toprule
%Graphs & Can & Be  & Fun & Too \\ 
%\midrule 
%$G_1$ & $1$  & $2$ & $3$ & $4$ \\ 
%$G_2$ & $-1$  & $-2$ & $-3$ & $-4$ \\ 
%$G_3$ & $0.1$  & $0.2$ & $0.3$ & $0.4$ \\ 
%\bottomrule
%\end{tabular}
%\label{tab}
%\end{table}

% A functor $F:\mathcal{C} \rightarrow \mathcal{D}$ is
% \begin{itemize}
%     \item \textbf{faithful} if each map $F_{A,B}:\mathcal{C}(A,B) \rightarrow \mathcal{D}(F\ A, F\ B)$ is injective;
%     \item \textbf{full} if each map $F_{A,B}:\mathcal{C}(A,B) \rightarrow \mathcal{D}(F\ A, F\ B)$ is surjective
% \end{itemize}

{\color{teal}
The following statements are \textbf{false} in general
{\color{red}
\begin{itemize}
    \item A functor $F :\mathcal{C}\rightarrow \mathcal{D}$ preserves limits iff $F$ preserves products and limits.
    \item A functor $F$ preserves products iff $F$ preserves binary products and terminal objects.
\end{itemize}
}
Only true if $\mathcal{C}$ is.

Functors preserve all commutative diagrams.
Functors preserve all cones.
Preserves limits $\implies$ sends universal cones to universal cones.

% https://q.uiver.app/#q=WzAsMyxbMCwwLCJcXG1hdGhjYWx7SX0iXSxbMSwwLCJcXG1hdGhjYWx7Q30iXSxbMiwwLCJcXG1hdGhjYWx7RH0iXSxbMCwxLCJKIiwyLHsiY3VydmUiOjF9XSxbMSwyLCJGIl0sWzAsMSwiXFxEZWx0YSBMIiwxLHsiY3VydmUiOi0xfV0sWzAsMiwiXFxEZWx0YSBGTCIsMCx7ImN1cnZlIjotM31dLFs1LDMsIiIsMCx7InNob3J0ZW4iOnsic291cmNlIjoyMCwidGFyZ2V0IjoyMH19XV0=
\[\begin{tikzcd}
	{\mathcal{I}} & {\mathcal{C}} & {\mathcal{D}}
	\arrow[""{name=0, anchor=center, inner sep=0}, "J"', curve={height=6pt}, from=1-1, to=1-2]
	\arrow["F", from=1-2, to=1-3]
	\arrow[""{name=1, anchor=center, inner sep=0}, "{\Delta L}"{description}, curve={height=-6pt}, from=1-1, to=1-2]
	\arrow["{\Delta FL}", curve={height=-18pt}, from=1-1, to=1-3]
	\arrow[shorten <=2pt, shorten >=2pt, Rightarrow, from=1, to=0]
\end{tikzcd}\]
}

\section*{Question 1}
\addtocounter{section}{1}
\begin{proposition}
    The covariant Hom-functor $\mathcal{C}(A,-):\mathcal{C}\rightarrow$ \textbf{Set} preserves all finite limits.
\end{proposition}

\begin{proof}
Let $D: \mathcal{J} \rightarrow \mathcal{C}$ be a diagram where $\mathcal{J}$ is a finite category.
Suppose the cone $(C, (c_1,c_2,...,c_N))$ to $D$ is a finite limit.
By being a cone, for any $f:I\rightarrow J$ in $\mathcal{J}$, we have
% https://q.uiver.app/#q=WzAsMyxbMCwwLCJEXFwgSSJdLFsyLDAsIkRcXCBKIl0sWzEsMSwiQyJdLFswLDEsIkRcXCBmIl0sWzIsMCwiY19JIl0sWzIsMSwiY19KIiwyXV0=
\[\begin{tikzcd}
	{D\ I} && {D\ J} \\
	& C
	\arrow["{D\ f}", from=1-1, to=1-3]
	\arrow["{c_I}", from=2-2, to=1-1]
	\arrow["{c_J}"', from=2-2, to=1-3]
\end{tikzcd}\]
\[D\ f\circ c_I = c_J\]

We show that $(\mathcal{C}(A,C), (\mathcal{C}(A,c_1), ,...,\mathcal{C}(A,c_N)))$ is a cone 
and a limit to $\mathcal{C}(A,-) \circ D: \mathcal{J}\rightarrow \textbf{Set}$.

% https://q.uiver.app/#q=WzAsOCxbMCwwLCJcXG1hdGhjYWx7Q30iXSxbMiwwLCJcXHRleHRiZntTZXR9Il0sWzAsMSwiQyJdLFswLDIsIkRcXCBJIl0sWzIsMSwiXFxtYXRoY2Fse0N9KEEsQykiXSxbMiwyLCJcXG1hdGhjYWx7Q30oQSwgRFxcIEkpIl0sWzAsMywiRFxcIEoiXSxbMiwzLCJcXG1hdGhjYWx7Q30oQSwgRFxcIEopIl0sWzAsMSwiXFxtYXRoY2Fse0N9KEEsLSkiXSxbMiwzLCJjX0kiLDJdLFs0LDUsIlxcbWF0aGNhbHtDfShBLGNfSSkiXSxbMiw0LCIiLDEseyJzdHlsZSI6eyJ0YWlsIjp7Im5hbWUiOiJtYXBzIHRvIn19fV0sWzMsNSwiIiwxLHsic3R5bGUiOnsidGFpbCI6eyJuYW1lIjoibWFwcyB0byJ9fX1dLFszLDYsIkRcXCBmIiwyXSxbNSw3LCJcXG1hdGhjYWx7Q30oQSwgRFxcIGYpIl0sWzYsNywiIiwxLHsic3R5bGUiOnsidGFpbCI6eyJuYW1lIjoibWFwcyB0byJ9fX1dLFs5LDEwLCIiLDIseyJzaG9ydGVuIjp7InNvdXJjZSI6MjAsInRhcmdldCI6MjB9LCJsZXZlbCI6MSwic3R5bGUiOnsidGFpbCI6eyJuYW1lIjoibWFwcyB0byJ9fX1dLFsxMywxNCwiIiwyLHsic2hvcnRlbiI6eyJzb3VyY2UiOjIwLCJ0YXJnZXQiOjIwfSwibGV2ZWwiOjEsInN0eWxlIjp7InRhaWwiOnsibmFtZSI6Im1hcHMgdG8ifX19XV0=
\[\begin{tikzcd}
	{\mathcal{C}} && {\textbf{Set}} \\
	C && {\mathcal{C}(A,C)} \\
	{D\ I} && {\mathcal{C}(A, D\ I)} \\
	{D\ J} && {\mathcal{C}(A, D\ J)}
	\arrow["{\mathcal{C}(A,-)}", from=1-1, to=1-3]
	\arrow[""{name=0, anchor=center, inner sep=0}, "{c_I}"', from=2-1, to=3-1]
	\arrow[""{name=1, anchor=center, inner sep=0}, "{\mathcal{C}(A,c_I)}", from=2-3, to=3-3]
	\arrow[maps to, from=2-1, to=2-3]
	\arrow[maps to, from=3-1, to=3-3]
	\arrow[""{name=2, anchor=center, inner sep=0}, "{D\ f}"', from=3-1, to=4-1]
	\arrow[""{name=3, anchor=center, inner sep=0}, "{\mathcal{C}(A, D\ f)}", from=3-3, to=4-3]
	\arrow[maps to, from=4-1, to=4-3]
	\arrow[shorten <=15pt, shorten >=15pt, maps to, from=0, to=1]
	\arrow[shorten <=15pt, shorten >=15pt, maps to, from=2, to=3]
\end{tikzcd}\]
Suppose we have an arbitrary $f:I\rightarrow J$ in $\mathcal{J}$,
so we get
\[\mathcal{C}(A,(D\ f)): \mathcal{C}(A,(D\ I)) \rightarrow\mathcal{C}(A,(D\ J))\]
Consider the morphisms $D\ I \xleftarrow[]{c_I}C\xrightarrow[]{c_J} D\ J$, so
\[\mathcal{C}(A,c_I):\mathcal{C}(A,C)\rightarrow \mathcal{C}(A, (D\ I)):= h\mapsto c_I \circ h\]
\[\mathcal{C}(A,c_J):\mathcal{C}(A,C)\rightarrow \mathcal{C}(A, (D\ J)):= h'\mapsto c_J \circ h'\]
% https://q.uiver.app/#q=WzAsMyxbMCwwLCJcXG1hdGhjYWx7Q30oQSwgRFxcIEkpIl0sWzIsMCwiXFxtYXRoY2Fse0N9KEEsIERcXCBJKSJdLFsxLDEsIlxcbWF0aGNhbHtDfShBLCBDKSJdLFswLDEsIlxcbWF0aGNhbHtDfShBLERcXCBmKSJdLFsyLDAsIlxcbWF0aGNhbHtDfShBLCBjX0kpIl0sWzIsMSwiXFxtYXRoY2Fse0N9KEEsIGNfSSkiLDJdXQ==
\[\begin{tikzcd}
	{\mathcal{C}(A, D\ I)} && {\mathcal{C}(A, D\ J)} \\
	& {\mathcal{C}(A, C)}
	\arrow["{\mathcal{C}(A,D\ f)}", from=1-1, to=1-3]
	\arrow["{\mathcal{C}(A, c_I)}", from=2-2, to=1-1]
	\arrow["{\mathcal{C}(A, c_J)}"', from=2-2, to=1-3]
\end{tikzcd}\]
Then the diagram commutes since
\begin{equation}
\begin{aligned}
    \mathcal{C}(A,(D\ f)) \circ \mathcal{C}(A,c_I) &= (h \mapsto (D\ f)\circ h) \circ (h' \mapsto c_I \circ h')\\
    &= h \mapsto (D\ f)\circ c_I \circ h\\
    &= h \mapsto c_J \circ h\\
    &= \mathcal{C}(A,c_J)
\end{aligned}
\end{equation}
So $(\mathcal{C}(A,C), (\mathcal{C}(A,c_1), ,...,\mathcal{C}(A,c_N)))$ is a cone to $\mathcal{C}(A,-) \circ D$.
    
To show the limit, consider $\mathcal{C}(A,C)$ in \textbf{Set} and let $(B,\gamma)$ 
be some other cone to  $\mathcal{C}(A,-) \circ D$, where $\gamma$ is
\[\{\gamma_I:B\rightarrow \mathcal{C}(A,(D\ I))\}_{I\in\text{Ob}(\mathcal{J})}\]
Since $(C, (c_1,c_2,...,c_N))$ is a limit to $D$, for any cone $(A, \delta)$ to $D$, 
there is a unique arrow $a:A\rightarrow C$ such that $\forall I\in\text{Ob}(\mathcal{J}),c_I \circ a = \delta_I$.

Note that $\forall b\in B,\gamma_I (b) \in \mathcal{C}(A,(D\ I))$.
Then for each $b\in B$, we have the cone 
\[(A,\{\gamma_I(b):A\rightarrow D\ I\}_{I\in\text{Ob}(\mathcal{J})})\]
From being the limit, there is a unique map $u_b:A\rightarrow C$ such that
% https://q.uiver.app/#q=WzAsMyxbMCwxLCJBIl0sWzEsMSwiQyJdLFsxLDAsIkRcXCBJIl0sWzAsMSwidV9iIiwyXSxbMCwyLCJcXGdhbW1hX0koYikiXSxbMSwyLCJjX0kiLDJdXQ==
\[\begin{tikzcd}
	& {D\ I} \\
	A & C
	\arrow["{u_b}"', from=2-1, to=2-2]
	\arrow["{\gamma_I(b)}", from=2-1, to=1-2]
	\arrow["{c_I}"', from=2-2, to=1-2]
\end{tikzcd}\]
\[\forall I\in \text{Ob}(\mathcal{J}), c_I\circ u_b = \gamma_I(b)\]

By defining the arrow $g:B\rightarrow \mathcal{C}(A,C)$ such that $\forall b\in B, g(b) = u_b$,
$g$ is the unique arrow from the cone $(B, \gamma)$ to 
$(\mathcal{C}(A,C), (\mathcal{C}(A,c_1), ,...,\mathcal{C}(A,c_N)))$, 
determined by the uniqueness of each $u_b$.
% https://q.uiver.app/#q=WzAsMyxbMCwxLCJCIl0sWzEsMSwiXFxtYXRoY2Fse0N9KEEsQykiXSxbMSwwLCJcXG1hdGhjYWx7Q30oQSxEXFwgSSkiXSxbMCwxLCJnIiwyXSxbMCwyLCJcXGdhbW1hX0kiXSxbMSwyLCJcXG1hdGhjYWx7Q30oQSxjX0kpIiwyXV0=
\[\begin{tikzcd}
	& {\mathcal{C}(A,D\ I)} \\
	B & {\mathcal{C}(A,C)}
	\arrow["g"', from=2-1, to=2-2]
	\arrow["{\gamma_I}", from=2-1, to=1-2]
	\arrow["{\mathcal{C}(A,c_I)}"', from=2-2, to=1-2]
\end{tikzcd}\]

Thus, $(\mathcal{C}(A,C), (\mathcal{C}(A,c_1), ,...,\mathcal{C}(A,c_N)))$ 
is the finite limit to $\mathcal{C}(A,-) \circ D$.
    
\end{proof}


\begin{corollary}
Since all representable functors are naturally isomorphic to the covariant Hom-functor at $A$,
$\mathcal{C}(A,-):C\rightarrow$ \textbf{Set} for some $A$, 
representable functors then preserve finite limits.
\end{corollary}

{\color{teal}
% https://q.uiver.app/#q=WzAsNixbMSwyLCJGTCJdLFswLDMsIkZKeCJdLFsyLDMsIkZKeSJdLFs0LDAsIlxcYnVsbGV0Il0sWzUsMSwiXFxidWxsZXQiXSxbMywxLCJcXGJ1bGxldCJdLFswLDFdLFswLDJdLFszLDRdLFszLDVdLFs1LDRdLFsxLDJdLFswLDMsIlxcY29uZyIsMV0sWzEsNSwiXFxjb25nIiwxXSxbMiw0LCJcXGNvbmciLDFdXQ==
\[\begin{tikzcd}
	&&&& \bullet \\
	&&& \bullet && \bullet \\
	& FL \\
	FJx && FJy
	\arrow[from=3-2, to=4-1]
	\arrow[from=3-2, to=4-3]
	\arrow[from=1-5, to=2-6]
	\arrow[from=1-5, to=2-4]
	\arrow[from=2-4, to=2-6]
	\arrow[from=4-1, to=4-3]
	\arrow["\cong"{description}, from=3-2, to=1-5]
	\arrow["\cong"{description}, from=4-1, to=2-4]
	\arrow["\cong"{description}, from=4-3, to=2-6]
\end{tikzcd}\]
\[\alpha : F \cong \mathcal{C}(A,-)\]

\[\mathcal{C}(W,A\times B) \cong \mathcal{C}(W,A)\times \mathcal{C}(W,B)\]
preserves binary products only if natural in $W$. 
Objectwise isomorphic functors not necessarily naturally isomorphic.
}

\clearpage
\section*{Question 2}
\addtocounter{section}{1}
\begin{proposition}
	The natural transformation $\Delta:$ Id $\rightarrow F$ 
	given by arrows $\Delta_X: X\rightarrow X\times X:=x\mapsto (x,x)$
	is the only natural transformation of the type Id $\rightarrow F$.
\end{proposition}
\begin{proof}
	$\Delta$ is a natural transformation since the following naturality square commutes,
	where $f: A\rightarrow B$ is an arbitrary arrow in \textbf{Set}, and $f(a)=b\in B$,
	with the mappings as shown, for every $a\in A, A\in \text{Ob}(\textbf{Set})$.
	% https://q.uiver.app/#q=WzAsMTAsWzIsMSwiYSJdLFszLDEsImIiXSxbMywyLCIoYixiKSJdLFsyLDIsIihhLGEpIl0sWzEsMCwiXFx0ZXh0e0lkIH1BPUEiXSxbNCwwLCJcXHRleHR7SWQgfUI9QiJdLFs0LDMsIkZcXCBCPUJcXHRpbWVzIEIiXSxbMSwzLCJGXFwgQT1BXFx0aW1lcyBBIl0sWzAsMCwiXFx0ZXh0e0lkfSJdLFswLDMsIkYiXSxbMCwxLCIiLDAseyJzdHlsZSI6eyJ0YWlsIjp7Im5hbWUiOiJtYXBzIHRvIn19fV0sWzEsMiwiIiwwLHsic3R5bGUiOnsidGFpbCI6eyJuYW1lIjoibWFwcyB0byJ9fX1dLFswLDMsIiIsMix7InN0eWxlIjp7InRhaWwiOnsibmFtZSI6Im1hcHMgdG8ifX19XSxbMywyLCIiLDIseyJzdHlsZSI6eyJ0YWlsIjp7Im5hbWUiOiJtYXBzIHRvIn19fV0sWzQsNSwiXFx0ZXh0e0lkIH0gZj1mIiwyXSxbNSw2LCJcXERlbHRhX0IiLDJdLFs0LDcsIlxcRGVsdGFfQSIsMl0sWzcsNiwiRlxcIGY9ZlxcdGltZXMgZiJdLFs4LDksIlxcRGVsdGEiLDJdXQ==
\[\begin{tikzcd}
	{\text{Id}} & {\text{Id }A=A} &&& {\text{Id }B=B} \\
	&& a & b \\
	&& {(a,a)} & {(b,b)} \\
	F & {F\ A=A\times A} &&& {F\ B=B\times B}
	\arrow[maps to, from=2-3, to=2-4]
	\arrow[maps to, from=2-4, to=3-4]
	\arrow[maps to, from=2-3, to=3-3]
	\arrow[maps to, from=3-3, to=3-4]
	\arrow["{\text{Id } f=f}"', from=1-2, to=1-5]
	\arrow["{\Delta_B}"', from=1-5, to=4-5]
	\arrow["{\Delta_A}"', from=1-2, to=4-2]
	\arrow["{F\ f=f\times f}", from=4-2, to=4-5]
	\arrow["\Delta"', from=1-1, to=4-1]
\end{tikzcd}\]
Suppose $\Gamma:\text{Id}\rightarrow F$ is another natural transformation.
Then for any function $f:\{*\}\rightarrow A$ on the singleton set where $f(*)=a\in A$, 
the following diagram commutes.
% https://q.uiver.app/#q=WzAsMTEsWzIsMSwiKiJdLFszLDEsImEiXSxbMywyLCIoYSxhKSJdLFsyLDIsIigqLCopIl0sWzEsMCwiXFx0ZXh0e0lkIH1cXHsqXFx9PVxceypcXH0iXSxbNCwwLCJcXHRleHR7SWQgfUE9QSJdLFs0LDMsIkZcXCBBPUFcXHRpbWVzIEEiXSxbMSwzLCJGXFwgXFx7KlxcfTo9XFx7KlxcfVxcdGltZXMgXFx7KlxcfSJdLFswLDAsIlxcdGV4dHtJZH0iXSxbMCwzLCJGIl0sWzEsNCwiPVxceygqLCopXFx9Il0sWzAsMSwiIiwwLHsic3R5bGUiOnsidGFpbCI6eyJuYW1lIjoibWFwcyB0byJ9fX1dLFsxLDIsIiIsMCx7InN0eWxlIjp7InRhaWwiOnsibmFtZSI6Im1hcHMgdG8ifX19XSxbMCwzLCIiLDIseyJzdHlsZSI6eyJ0YWlsIjp7Im5hbWUiOiJtYXBzIHRvIn19fV0sWzMsMiwiIiwyLHsic3R5bGUiOnsidGFpbCI6eyJuYW1lIjoibWFwcyB0byJ9fX1dLFs0LDUsIlxcdGV4dHtJZCB9IGY9ZiJdLFs1LDYsIlxcR2FtbWFfQSIsMl0sWzQsNywiXFxHYW1tYV97XFx7KlxcfX0iLDJdLFs3LDYsIkZcXCBmPWZcXHRpbWVzIGYiLDJdLFs4LDksIlxcR2FtbWEiLDJdXQ==
\[\begin{tikzcd}
	{\text{Id}} & {\text{Id }\{*\}=\{*\}} &&& {\text{Id }A=A} \\
	&& {*} & a \\
	&& {(*,*)} & {(a,a)} \\
	F & {F\ \{*\}:=\{*\}\times \{*\}} &&& {F\ A=A\times A} \\
	& {=\{(*,*)\}}
	\arrow[maps to, from=2-3, to=2-4]
	\arrow[maps to, from=2-4, to=3-4]
	\arrow[maps to, from=2-3, to=3-3]
	\arrow[maps to, from=3-3, to=3-4]
	\arrow["{\text{Id } f=f}", from=1-2, to=1-5]
	\arrow["{\Gamma_A}"', from=1-5, to=4-5]
	\arrow["{\Gamma_{\{*\}}}"', from=1-2, to=4-2]
	\arrow["{F\ f=f\times f}"', from=4-2, to=4-5]
	\arrow["\Gamma"', from=1-1, to=4-1]
\end{tikzcd}\]
In other words, we must have
\[\forall A\in \text{Ob}(\textbf{Set}): \Gamma_A (a)=(f\times f) (*,*) = (f(*),f(*))= (a,a)\]
Then $\forall A, \Gamma_A = \Delta_A$, so $\Delta$ is the unique natural transformation.
\end{proof}

\clearpage
\section*{Question 3}
\addtocounter{section}{1}
Define 
\[\texttt{eval}:\mathcal{D}^{\mathcal{C}}\times \mathcal{C} \rightarrow \mathcal{D}\]
\[\texttt{eval}::(F:\mathcal{C}\rightarrow\mathcal{D},C) \mapsto F\ C, \quad
(\Delta:F\rightarrow G, h:C\rightarrow C') \mapsto (k : F\ C \rightarrow G\ C')\]
{\color{teal} $k$ comes from the naturality square where for $f:C\rightarrow C'$,
\[G\ f \circ \Delta_C = \Delta_{C'} \circ F\ f\]}
% https://q.uiver.app/#q=WzAsNixbMCwwLCJcXG1hdGhjYWx7RH1ee1xcbWF0aGNhbHtDfX1cXHRpbWVzIFxcbWF0aGNhbHtDfSJdLFsyLDAsIlxcbWF0aGNhbHtEfSJdLFswLDEsIihGOlxcbWF0aGNhbHtDfVxccmlnaHRhcnJvd1xcbWF0aGNhbHtEfSxDKSJdLFsyLDEsIkZcXCBDIl0sWzAsMiwiKEY6XFxtYXRoY2Fse0N9XFxyaWdodGFycm93XFxtYXRoY2Fse0R9LEMnKSJdLFsyLDIsIkdcXCBDJyJdLFswLDEsIlxcdGV4dHtldmFsfSJdLFsyLDMsIiIsMCx7InN0eWxlIjp7InRhaWwiOnsibmFtZSI6Im1hcHMgdG8ifX19XSxbMiw0LCIoXFxEZWx0YTpGXFxyaWdodGFycm93IEcsaDpDXFxyaWdodGFycm93IEMnKSIsMl0sWzMsNSwiXFxHYW1tYTpGXFwgQ1xccmlnaHRhcnJvdyBHXFwgQyciXSxbNCw1LCIiLDIseyJzdHlsZSI6eyJ0YWlsIjp7Im5hbWUiOiJtYXBzIHRvIn19fV0sWzgsOSwiIiwyLHsic2hvcnRlbiI6eyJzb3VyY2UiOjIwLCJ0YXJnZXQiOjIwfSwibGV2ZWwiOjEsInN0eWxlIjp7InRhaWwiOnsibmFtZSI6Im1hcHMgdG8ifX19XV0=
\[\begin{tikzcd}
	{\mathcal{D}^{\mathcal{C}}\times \mathcal{C}} && {\mathcal{D}} \\
	{(F:\mathcal{C}\rightarrow\mathcal{D},C)} && {F\ C} \\
	{(F:\mathcal{C}\rightarrow\mathcal{D},C')} && {G\ C'}
	\arrow["{\texttt{eval}}", from=1-1, to=1-3]
	\arrow[maps to, from=2-1, to=2-3]
	\arrow[""{name=0, anchor=center, inner sep=0}, "{(\Delta:F\rightarrow G,h:C\rightarrow C')}"', from=2-1, to=3-1]
	\arrow[""{name=1, anchor=center, inner sep=0}, "{k:F\ C\rightarrow G\ C'}", from=2-3, to=3-3]
	\arrow[maps to, from=3-1, to=3-3]
	\arrow[shorten <=17pt, shorten >=17pt, maps to, from=0, to=1]
\end{tikzcd}\]
\begin{lemma}
	$\texttt{eval}:\mathcal{D}^{\mathcal{C}}\times \mathcal{C} \rightarrow \mathcal{D}$ is a functor.
	\label{lem:1}
\end{lemma}
\begin{proof}
We use the lemma to show that \texttt{eval} is a functor.
Fix an object $F:\mathcal{C} \rightarrow \mathcal{D}$ in $\mathcal{D}^{\mathcal{C}}$ so 
\[\texttt{eval}(F,-):\mathcal{C}\rightarrow\mathcal{D}\]
\[\texttt{eval}(F,-):: C\mapsto \texttt{eval}(F,C), \quad f \mapsto F\ f\]
Clearly each object $C$ in $\mathcal{C}$ has a mapping to $F\ C$ in $\mathcal{D}$.
For arrow $f:A\rightarrow B,g:B\rightarrow C$ in $\mathcal{C}$,
\begin{equation}
\begin{aligned}
	\texttt{eval}(F,g\circ f)&=F(g\circ f)\\
	&=F\ g \circ F\ f \quad\text{(by functoriality)}\\
	&= \texttt{eval}(F,g) \circ \texttt{eval}(F,g)
\end{aligned}
\end{equation}
\[\texttt{eval}(F,1_A)=F\ 1_A = 1_{F\ A} = 1_{\texttt{eval}(F,A)}\]
So $\texttt{eval}(F,-)$ satisfies functoriality.

Now fix an object $C$ in $\mathcal{C}$ so
\[\texttt{eval}(-,C):\mathcal{D}^{\mathcal{C}}\rightarrow\mathcal{D}\]
\[\texttt{eval}(-,C):: F \mapsto \texttt{eval}(F,C), \quad t \mapsto t_C\]
Each functor $F:\mathcal{C}\rightarrow \mathcal{D}$ in $\mathcal{D}^{\mathcal{C}}$
has a mapping to $F\ C$ in $\mathcal{D}$.
For each natural transformation $\Delta:F\rightarrow G, \Gamma:G\rightarrow H$ in $\mathcal{D}^{\mathcal{C}}$,
% https://q.uiver.app/#q=WzAsMyxbMCwwLCJGXFwgQyJdLFsxLDAsIkdcXCBDIl0sWzIsMCwiSFxcIEMiXSxbMCwxLCJcXERlbHRhX0MiXSxbMSwyLCJcXEdhbW1hX0MiXSxbMCwyLCIoXFxHYW1tYVxcY2lyY1xcRGVsdGEpX0MiLDAseyJjdXJ2ZSI6LTN9XV0=
\[\begin{tikzcd}
	{F\ C} & {G\ C} & {H\ C}
	\arrow["{\Delta_C}", from=1-1, to=1-2]
	\arrow["{\Gamma_C}", from=1-2, to=1-3]
	\arrow["{(\Gamma\circ\Delta)_C}", curve={height=-18pt}, from=1-1, to=1-3]
\end{tikzcd}\]
\begin{equation}
\begin{aligned}
	\texttt{eval}(\Gamma\circ\Delta, C) &= (\Gamma\circ\Delta)_C:F\ C\rightarrow H\ C\\
	&= (\Gamma_C:G\ C\rightarrow H\ C) \circ (\Delta_C : F\ C \rightarrow H\ C)\\
	&= \texttt{eval}(\Gamma, C) \circ \texttt{eval}(\Delta, C)
\end{aligned}
\end{equation}
\[\texttt{eval}(1_F, C) = 1_{F\ C} = 1_{\texttt{eval}(F,C)}\]
So $\texttt{eval}(-,C)$ is a functor, and the first part of the lemma is satisfied.

Suppose we have a natural transformation $\Delta: F\rightarrow G$ in $\mathcal{D}^{\mathcal{C}}$,
and an arrow $f:C\rightarrow C'$ in $\mathcal{C}$.
% https://q.uiver.app/#q=WzAsNCxbMCwwLCJcXHRleHR7ZXZhbH0oRixDKSJdLFsxLDAsIlxcdGV4dHtldmFsfShGLEMnKSJdLFsxLDEsIlxcdGV4dHtldmFsfShHLEMnKSJdLFswLDEsIlxcdGV4dHtldmFsfShHLEMpIl0sWzAsMSwiXFx0ZXh0e2V2YWx9KEYsZikiXSxbMSwyLCJcXHRleHR7ZXZhbH0oXFxEZWx0YSxDJykiXSxbMCwzLCJcXHRleHR7ZXZhbH0oXFxEZWx0YSxDKSIsMl0sWzMsMiwiXFx0ZXh0e2V2YWx9KEcsZikiXV0=
\[\begin{tikzcd}
	{\texttt{eval}(F,C)} & {\texttt{eval}(F,C')} \\
	{\texttt{eval}(G,C)} & {\texttt{eval}(G,C')}
	\arrow["{\texttt{eval}(F,f)}", from=1-1, to=1-2]
	\arrow["{\texttt{eval}(\Delta,C')}", from=1-2, to=2-2]
	\arrow["{\texttt{eval}(\Delta,C)}"', from=1-1, to=2-1]
	\arrow["{\texttt{eval}(G,f)}", from=2-1, to=2-2]
\end{tikzcd}\]
Then the diagram commutes since the naturality square commutes with $\Delta$ being a natural transformation.
% https://q.uiver.app/#q=WzAsNCxbMCwwLCJGXFwgQyJdLFsxLDAsIkZcXCBDJyJdLFsxLDEsIkdcXCBDJyJdLFswLDEsIkdcXCBDIl0sWzAsMSwiRlxcIGYiXSxbMSwyLCJcXERlbHRhX3tDJ30iXSxbMCwzLCJcXERlbHRhX0MiLDJdLFszLDIsIkdcXCBmIl1d
\[\begin{tikzcd}
	{F\ C} & {F\ C'} \\
	{G\ C} & {G\ C'}
	\arrow["{F\ f}", from=1-1, to=1-2]
	\arrow["{\Delta_{C'}}", from=1-2, to=2-2]
	\arrow["{\Delta_C}"', from=1-1, to=2-1]
	\arrow["{G\ f}", from=2-1, to=2-2]
\end{tikzcd}\]
\begin{equation}
\begin{aligned}
	\texttt{eval}(F,f) \circ \texttt{eval}(\Delta,C')&=(F\ f) \circ \Delta_C'\\
	&= (G\ f) \circ \Delta_C \quad\text{(by naturality)}\\
	&= \texttt{eval}(G,f) \circ \texttt{eval}(\Delta,C)
\end{aligned}
\end{equation}
The second part of the lemma is satisfied, and 
$\texttt{eval}:\mathcal{D}^{\mathcal{C}}\times \mathcal{C} \rightarrow \mathcal{D}$
is a functor.
\end{proof}

Let $G: \mathcal{E} \times \mathcal{C} \rightarrow \mathcal{D}$ be an arbitrary functor.
Define 
\[\tilde{G} : \mathcal{E} \rightarrow \mathcal{D}^{\mathcal{C}}\]
\[\tilde{G} :: A \mapsto G(A, -),\quad (f:A\rightarrow B) \mapsto (\Delta: G(A,-) \rightarrow G(B,-))\]
where $G(A,-):\mathcal{C} \rightarrow \mathcal{D}$ is a functor for each $A$ in $\mathcal{E}$.
\begin{lemma}
For a given $G: \mathcal{E} \times \mathcal{C} \rightarrow \mathcal{D}$,
$\tilde{G} : \mathcal{E} \rightarrow \mathcal{D}^{\mathcal{C}}$ is the unique functor
where the following diagram commutes.
% https://q.uiver.app/#q=WzAsMyxbMCwxLCJcXG1hdGhjYWx7RX1cXHRpbWVzIFxcbWF0aGNhbHtDfSJdLFswLDAsIlxcbWF0aGNhbHtEfV57XFxtYXRoY2Fse0N9fVxcdGltZXMgXFxtYXRoY2Fse0N9Il0sWzEsMCwiXFxtYXRoY2Fse0R9Il0sWzAsMSwiXFx0aWxkZXtHfVxcdGltZXMgMV97XFxtYXRoY2Fse0N9fSJdLFswLDIsIkciLDJdLFsxLDIsIlxcdGV4dHtldmFsfSJdXQ==
\[\begin{tikzcd}
	{\mathcal{D}^{\mathcal{C}}\times \mathcal{C}} & {\mathcal{D}} \\
	{\mathcal{E}\times \mathcal{C}}
	\arrow["{\tilde{G}\times 1_{\mathcal{C}}}", from=2-1, to=1-1]
	\arrow["G"', from=2-1, to=1-2]
	\arrow["{\texttt{eval}}", from=1-1, to=1-2]
\end{tikzcd}\]
\label{lem:2}
\end{lemma}

\begin{proof}
Any object $A$ in $\mathcal{E}$ will get mapped to $G(A,-)$ in $\mathcal{D}^{\mathcal{C}}$.
Suppose $A\xrightarrow[]{f}A'\xrightarrow[]{g}A''$ in $\mathcal{E}$. Then
\begin{equation}
\begin{aligned}
	\forall C\in\text{Ob}(\mathcal{C}), ((\tilde{G}\ g) \circ (\tilde{G}\ f))(C)
	&= (\tilde{G}\ g)(C) \circ (\tilde{G}\ f)(C)\\
	&= (\Delta'_C:G(A',C)\rightarrow G(A'',C)) \circ (\Delta_C: G(A,C)\rightarrow G(A',C))\\
	&= (\Delta'\circ \Delta)_C: G(A,C) \rightarrow G(A'',C)\\
	&= \tilde{G} (g\circ f)(C)
\end{aligned}
\end{equation}
\[\tilde{G}\ 1_A = 1_{G(A,-)}=1_{\tilde{G}\ A}\]
So $\tilde{G} : \mathcal{E} \rightarrow \mathcal{D}^{\mathcal{C}}$ is a functor.


Let $(E,C)$ be an object in $\mathcal{E}\times \mathcal{C}$.
\[\texttt{eval}\circ (\tilde{G}\times 1_{\mathcal{C}})(E,C)=\texttt{eval}(G(E,-),C) = G(E,C)\]
Let $(f,g)$ be an arrow in $\mathcal{E}\times \mathcal{C}$.
\begin{equation}
\begin{aligned}
	\texttt{eval}\circ (\tilde{G}\times 1_{\mathcal{C}})(f:A\rightarrow A', g:C\rightarrow C')
	&=\texttt{eval}(\Delta:G(A,-)\rightarrow G(A',-),g:C\rightarrow C')\\
	&=k: G(A,C)\rightarrow G(A', C')\\
	&= G(f:A\rightarrow A',g:C\rightarrow C')
\end{aligned}
\end{equation}
So $\texttt{eval}\circ (\tilde{G}\times 1_{\mathcal{C}}) = G$.

To show uniqueness, suppose there is some other functor $H: \mathcal{E} \rightarrow \mathcal{D}^{\mathcal{C}}$
such that $\texttt{eval}\circ (H\times 1_{\mathcal{C}}) = G$. 
Then for each object $(E,C)$ in $\mathcal{E}\times \mathcal{C}$,
\[\texttt{eval}\circ (H\times 1_{\mathcal{C}})(E,C)=\texttt{eval}(H(E),C) = H(E)(C)= G(E,C)\]
So $\forall (E,C)\in\text{Ob}(\mathcal{E}\times\mathcal{C}), H(E)(C)=G(E,C)$.
Then for each arrow $(f,g)$ in $\mathcal{E}\times \mathcal{C}$,
\begin{equation}
\begin{aligned}
	\texttt{eval}\circ (H\times 1_{\mathcal{C}})(f:A\rightarrow A', g:C\rightarrow C')
	&=\texttt{eval}(H\ f: H\ A\ \rightarrow H\ A',g:C\rightarrow C')\\
	&=k: H\ A\ C\rightarrow H\ A'\ C'\\
	&= G(f:A\rightarrow A',g:C\rightarrow C')\\
	&= k: G(A, C) \rightarrow G(A', C')
\end{aligned}
\end{equation}
So $\forall (f,g)\in\text{Ob}(\mathcal{E}\times\mathcal{C}), H(f)(g)=G(f,g)$.
Thus, $H=\tilde{G}$ is unique. 
\end{proof}
{\color{teal}
Exponentials are right adjoint to products.}

\begin{proposition}
	Func$(\mathcal{C},\mathcal{D})$ is an exponential object in \textbf{Cat}.
\end{proposition}
\begin{proof}
By Lemmas \ref{lem:1}, \ref{lem:2}, there is an object $\mathcal{D}^{\mathcal{C}}$ 
and morphism \texttt{eval}$:\mathcal{D}^{\mathcal{C}}\times \mathcal{C} \rightarrow \mathcal{D}$,
where for every $G: \mathcal{E} \times \mathcal{C} \rightarrow \mathcal{D}$,
there is a unique morphism $\tilde{G} : \mathcal{E} \rightarrow \mathcal{D}^{\mathcal{C}}$
such that $\texttt{eval}\circ (\tilde{G}\times 1_{\mathcal{C}}) = G$.
Func$(\mathcal{C},\mathcal{D})$ is an exponential object in \textbf{Cat}.
\end{proof}

\clearpage
\section*{Question 4}
\addtocounter{section}{1}
\subsection*{(a)}
\begin{proposition}
	$a\Rightarrow b$ is an exponential.
\end{proposition}
\begin{proof}
\[\texttt{eval}: (a\Rightarrow b) \wedge a \leq b\]
Let $f: c\wedge a \leq b$ be an arrow.
We want to show the unique arrow $\tilde{f}: c\leq (a\Rightarrow b)$ such that
\[\texttt{eval} \circ (\tilde{f}\times 1_a) = f\]
\begin{equation}
\begin{aligned}
	\texttt{eval} \circ (\tilde{f}\times 1_a) &= \texttt{eval} \circ (c\leq (a\Rightarrow b) \wedge a \leq a)\\
	&= ((a\Rightarrow b) \wedge a \leq b) \circ (c \wedge a \leq (a\Rightarrow b) \wedge a)\\
	&= c \wedge a \leq b\\
	&= f
\end{aligned}
\end{equation}
Since $B$ is a poset, there is at most 1 arrow between objects - $\tilde{f}$ is a unique arrow.

$a\Rightarrow b$ is an exponential.
	
\end{proof}
{\color{teal}
Finite boolean algebra $B$ is a poset $\implies$ thing category, which is closed (has exponentials)

Cartesian closed categories (CCC).

$\tilde{G}$ corresponds to $\lambda$ abstraction in functional programming.
% https://q.uiver.app/#q=WzAsNSxbMSwwLCJBXFxSaWdodGFycm93IEJcXHRpbWVzIEEiXSxbMiwwLCJCIl0sWzEsMSwiXFxHYW1tYSBcXHRpbWVzIEEiXSxbMCwxLCJcXEdhbW1hIl0sWzAsMCwiQVxcUmlnaHRhcnJvdyBCIl0sWzAsMSwiYXBwIl0sWzIsMSwidCIsMl0sWzIsMCwiXFxsYW1iZGEgdFxcdGltZXMgMV9BIl0sWzMsNCwiXFxsYW1iZGEgdCJdXQ==
\[\begin{tikzcd}
	{A\Rightarrow B} & {A\Rightarrow B\times A} & B \\
	\Gamma & {\Gamma \times A}
	\arrow["app", from=1-2, to=1-3]
	\arrow["t"', from=2-2, to=1-3]
	\arrow["{\lambda t\times 1_A}", from=2-2, to=1-2]
	\arrow["{\lambda t}", from=2-1, to=1-1]
\end{tikzcd}\]
\[(\lambda x. t)(a) = t[a/x]\]


Show
\[a\wedge b \leq c \iff b \leq a \implies c\]
\[\texttt{eval} : (\neg a \vee b) \wedge a = (\neg a \wedge a) \vee (b \wedge a) = b \wedge a\]
\begin{proof}
\[a\times b \rightarrow c\]
\begin{equation}
\begin{aligned}
    a \wedge b \leq c &\iff a \wedge b \leq \neg\neg c\\
    &\iff a \wedge b \wedge \neg c = \bot\\
    &\iff b \leq \neg(a \wedge \neg c)\\
    &\iff b \leq \neg a \vee c\\
    &\iff b \leq a \implies c
\end{aligned}
\end{equation}
\[b \rightarrow c^a\]
\end{proof}
}

\subsection*{(b)}
\[\bot=\emptyset\]
\[\top=X\]

\subsection*{(c)}
The atoms are the elements of the set $X$ as sets. i.e. the set $\{x\}$ where $x\in X$.

\clearpage
\subsection*{(d)}
\begin{proposition}
	$t:\text{Id}\rightarrow At\circ \mathcal{P}^B$ is a natural isomorphism.
\end{proposition}
% https://q.uiver.app/#q=WzAsMTAsWzIsMSwiYSJdLFszLDEsImYoYSkiXSxbMywyLCJcXHtmKGEpXFx9Il0sWzIsMiwiXFx7YVxcfSJdLFsxLDAsIkEiXSxbNCwwLCJCIl0sWzQsMywiXFx7XFx7YlxcfXxiXFxpbiBCXFx9Il0sWzEsMywiXFx7XFx7YVxcfXxhXFxpbiBBXFx9Il0sWzAsMCwiXFx0ZXh0e0lkfSJdLFswLDMsIkF0XFxjaXJjXFxtYXRoY2Fse1B9XkIiXSxbMCwxLCIiLDAseyJzdHlsZSI6eyJ0YWlsIjp7Im5hbWUiOiJtYXBzIHRvIn19fV0sWzEsMiwiIiwwLHsic3R5bGUiOnsidGFpbCI6eyJuYW1lIjoibWFwcyB0byJ9fX1dLFswLDMsIiIsMix7InN0eWxlIjp7InRhaWwiOnsibmFtZSI6Im1hcHMgdG8ifX19XSxbMywyLCIiLDIseyJzdHlsZSI6eyJ0YWlsIjp7Im5hbWUiOiJtYXBzIHRvIn19fV0sWzQsNSwiZiIsMl0sWzUsNiwidF9CIiwyLHsib2Zmc2V0IjoxfV0sWzQsNywidF9BIiwyLHsib2Zmc2V0IjoxfV0sWzcsNiwiXFx7YVxcfVxcbWFwc3RvIFxce2YoYSlcXH0iXSxbOCw5LCJ0IiwyLHsib2Zmc2V0IjoxfV0sWzksOCwidF57LTF9IiwyLHsib2Zmc2V0IjoyfV0sWzcsNCwidF9BXnstMX0iLDIseyJvZmZzZXQiOjF9XSxbNiw1LCJ0X0Jeey0xfSIsMix7Im9mZnNldCI6MX1dXQ==
\[\begin{tikzcd}
	{\text{Id}} & A &&& B \\
	&& a & {f(a)} \\
	&& {\{a\}} & {\{f(a)\}} \\
	{At\circ\mathcal{P}^B} & {\{\{a\}|a\in A\}} &&& {\{\{b\}|b\in B\}}
	\arrow[maps to, from=2-3, to=2-4]
	\arrow[maps to, from=2-4, to=3-4]
	\arrow[maps to, from=2-3, to=3-3]
	\arrow[maps to, from=3-3, to=3-4]
	\arrow["f"', from=1-2, to=1-5]
	\arrow["{t_B}"', shift right, from=1-5, to=4-5]
	\arrow["{t_A}"', shift right, from=1-2, to=4-2]
	\arrow["{\{a\}\mapsto \{f(a)\}}", from=4-2, to=4-5]
	\arrow["t"', shift right, from=1-1, to=4-1]
	\arrow["{t^{-1}}"', shift right=2, from=4-1, to=1-1]
	\arrow["{t_A^{-1}}"', shift right, from=4-2, to=1-2]
	\arrow["{t_B^{-1}}"', shift right, from=4-5, to=1-5]
\end{tikzcd}\]
\begin{proof}
To prove the naturality square commutes, we show that for all finite sets $A,B$,
\[((At\circ \mathcal{P}^B)\ f) \circ t_A = t_B \circ f\]

% https://q.uiver.app/#q=WzAsOSxbMSwxLCJcXG1hdGhjYWx7UH0oQSkiXSxbMiwxLCJcXHtcXHthXFx9fGFcXGluIEFcXH0iXSxbMCwyLCJCIl0sWzEsMiwiXFxtYXRoY2Fse1B9KEIpIl0sWzAsMCwiXFx0ZXh0YmZ7RmluU2V0fSJdLFsxLDAsIlxcdGV4dGJme0ZpbkJvb2x9XntvcH0iXSxbMiwwLCJcXHRleHRiZntGaW5TZXR9Il0sWzAsMSwiQSJdLFsyLDIsIlxce1xce2JcXH18YlxcaW4gQlxcfSJdLFswLDEsIiIsMCx7InN0eWxlIjp7InRhaWwiOnsibmFtZSI6Im1hcHMgdG8ifX19XSxbNCw1LCJcXG1hdGhjYWx7UH1eQiJdLFs1LDYsIkF0Il0sWzcsMCwiIiwwLHsic3R5bGUiOnsidGFpbCI6eyJuYW1lIjoibWFwcyB0byJ9fX1dLFs3LDIsImYiLDJdLFszLDAsImgoVCk9XFx7YVxcaW4gQXxmKGEpXFxpbiBUXFx9IiwxXSxbMSw4LCJBdChoKSJdLFsyLDMsIiIsMix7InN0eWxlIjp7InRhaWwiOnsibmFtZSI6Im1hcHMgdG8ifX19XSxbMyw4LCIiLDAseyJzdHlsZSI6eyJ0YWlsIjp7Im5hbWUiOiJtYXBzIHRvIn19fV1d
\[\begin{tikzcd}
	{\textbf{FinSet}} & {\textbf{FinBool}^{op}} & {\textbf{FinSet}} \\
	A & {\mathcal{P}(A)} & {\{\{a\}|a\in A\}} \\
	B & {\mathcal{P}(B)} & {\{\{b\}|b\in B\}}
	\arrow[maps to, from=2-2, to=2-3]
	\arrow["{\mathcal{P}^B}", from=1-1, to=1-2]
	\arrow["At", from=1-2, to=1-3]
	\arrow[maps to, from=2-1, to=2-2]
	\arrow["f"', from=2-1, to=3-1]
	\arrow["{h(T)=\{a\in A|f(a)\in T\}}"{description}, from=3-2, to=2-2]
	\arrow["{At(h)}", from=2-3, to=3-3]
	\arrow[maps to, from=3-1, to=3-2]
	\arrow[maps to, from=3-2, to=3-3]
\end{tikzcd}\]
Note that by definition of the functors $\mathcal{P}^B$ and $At$,
where $\mathcal{P}^B\ f = h:\mathcal{P}(B) \rightarrow \mathcal{P}(A)$,
\begin{equation}
\begin{aligned}
	\forall a\in A, &(\forall T\in \mathcal{P}(B), f(a)\in T \iff a\in h(T) \iff \{a\} \subseteq h(T))\\
	&\iff (\exists! \{b'\}\in \mathcal{P}(B): \{b'\} \subseteq T \iff b' \in T)
	\label{eq:10}
\end{aligned}
\end{equation}
In other words, for all elements $a\in A$, let its corresponding atom $\{a\}\subseteq h(T)$ 
for all subsets $T\subseteq B$ where $f(a)\in T$, so that $\{b'\} \subseteq B$ is the unique atom 
that satisfies the property of a Boolean algebra homomorphism.

It must be that $b'=f(a)$ since $\{f(a)\}\subseteq B$ is one such atom that satisfies the property,
and by uniqueness.
So $\forall a\in A,At(h)(\{a\})=\{f(a)\}\subseteq B$.
 $t$ is a natural transformation.

To show the natural isomorphism, we can define the inverse for all finite sets $A$,
\[t_A^{-1} : \{\{a\}|a\in A\} \rightarrow A\]
\[t_A^{-1} :: \{a\} \mapsto a, \quad (k: \{a\} \mapsto \{b\}) \mapsto (f(a) = b)\]
so $t_A^{-1}\circ t_A = 1_{\text{Id}\ A},t_A\circ t_A^{-1} = 1_{At\circ \mathcal{P}^B\ A} $.
We have for all $A, t_A$ is an isomorphism, so $t$ is a natural isomorphism.
	
\end{proof}

\section*{Optional}
\addtocounter{section}{1}
\subsection*{(a)}
Suppose we have a function $f:A\rightarrow B$.
Then $[f]\circ \texttt{flatten}_A=\texttt{flatten}_B\circ (\texttt{Tree}\ f)$
since each element in the underlying tree will get mapped to the same element in the resulting list.

\subsection*{(b)}
If we have the following where \texttt{undef1 . id = undef2},
\begin{lstlisting}[language=Haskell]
undef1 = undefined :: a -> b
undef2 = \_ -> undefined
seq undef1 () = undefined
seq undef2 () = ()
\end{lstlisting}
Different values are given with \texttt{seq}. 
The identity axiom is not satisfied so \textbf{Hask} is not a category.


\end{document}
